\bookchapter{제2장: 디지털 방을 정리하세요 - 파일 관리}

\chapterintro{스마트폰을 오래 사용하면 어지러워지며, 좋은 습관은 파일 관리에서 시작됩니다. 방을 정기적으로 정리해야 하듯이, 디지털 공간도 관리해야 합니다. 정돈된 스마트폰은 물건을 찾기 쉬울 뿐만 아니라, 실행도 더 매끄러움.1. 3계층 저장 체계: 당신의 디지털 창고\textbf{기기 저장}: 핵심 영역, 속도가 빠르며, 시스템과 앱이 모두 여기에 있음. 이것은 스마트폰에서 가장 귀중한 공간이며, 정밀하게 계획해야 함.\textbf{SD 카드 (지원되는 경우)}: 확장 옷장, 속도가 약간 느리며, 영화, 사진 등 대용량 파일에 적합. 참고: 현재 많은 플래그십 기기가 SD 카드를 지원하지 않음.\textbf{클라우드 저장}: 금고+동기화 네트워크 드라이브. 분실 방지에, 여러 기기 간 동기화도 가능. iCloud, Baidu Netdisk, OneDrive가 모두 일반적인 선택.2. “파일 관리자” 이해안드로이드에는 보통 [파일] 또는 [파일 관리] 앱이 있으며, iOS는 자체 [파일] 앱이 있음. 이것들은 스마트폰의 모든 파일을 볼 수 있는 “리소스 탐색기”임.일반적인 폴더 의미:\textbf{DCIM}: 카메라로 촬영한 사진과 영상.\textbf{Pictures}: 스크린샷, 저장된 이미지 등 기타 이미지.\textbf{Download}: 브라우저에서 다운로드한 파일.\textbf{Music}:음악파일.\textbf{Movies}:영상파일.\textbf{Documents}: 문서 파일.\textbf{Android}: 안드로이드 시스템 폴더, 각 앱의 데이터가 들어 있으며, 일반적으로 수동으로 삭제하지 않음.3. 파일 시스템 구조: 깊이 있는 이해안드로이드는 Linux 기반이며, 파일 시스템은 일정한 계층 구조를 가짐. 일반 사용자가 깊이 이해할 필요는 없지만, 몇 가지 기본 개념을 알면 도움이 됨.저장 파티션:\textbf{/system}: 시스템 파티션, 운영체제 파일 저장. 일반 사용자는 액세스 권한이 없으며, root 권한 필요.\textbf{/data}: 데이터 파티션, 사용자 데이터와 앱 데이터 저장. 이것은 일상적으로 사용하는 공간.\textbf{/sdcard 또는 /storage/emulated/0}: 사용자가 액세스할 수 있는 저장 공간, 사진, 다운로드 파일이 모두 여기에 있음.앱 데이터 저장:\textbf{내부 저장}: 앱의 개인 공간, 다른 앱은 액세스 불가. 앱 제거 시 삭제됨.\textbf{외부 저장}: 사용자가 액세스할 수 있는 영역에 앱이 저장하는 파일로, 사진, 다운로드 콘텐츠 등. 앱 제거 시 반드시 삭제되는 것은 아님.4. 데이터 백업과 복구: 데이터에 보험 들기데이터는 무가치하며, 백업은 반드시 길러야 할 습관입니다.백업 전략:\textbf{클라우드 백업}: 자동, 편리하지만 네트워크와 클라우드 공간 필요. iCloud, 화웨이 클라우드, 샤오미 클라우드 등 모두 지원.\textbf{로컬 백업}: 데이터 케이블로 컴퓨터 연결, 공식 도구로 백업. 속도 빠름, 공간 큼, 단 수동 조작 필요.\textbf{이중 백업}: 클라우드 백업+로컬 백업, 가장 안전.무엇을 백업할까?\textbf{사진과 영상}: 공간 가장 많이 차지, 가장 소중함.\textbf{연락처}: 클라우드 계정에 저장하면 가장 편리.\textbf{위챗 채팅 기록}: 위챗 자체 백업 기능 사용.\textbf{중요 문서}: 클라우드 드라이브나 컴퓨터에 저장.5. 파일 보안: 프라이버시 보호\textbf{민감 파일 암호화}: 많은 스마트폰에 “보안 파일함” 또는 “프라이빗 공간” 기능이 있어, 민감 파일을 넣을 수 있으며, 비밀번호나 지문이 있어야 접근 가능.\textbf{삭제는 파괴가 아님}: 일반 삭제는 공간 덮어쓰기 표시만 할 뿐, 데이터는 여전히 존재. 철저히 삭제하려면 전용 보안 삭제 도구를 사용하거나, 여러 번 덮어쓰기 필요.\textbf{클라우드 보안}: 평판 좋은 클라우드 서비스 공급자를 선택하고, 2단계 인증을 켜고, 로그인 기기를 정기적으로 확인.6. 공간이 부족한가? 대청소를 시작하세요\textbf{캐시 정리}: 앱의 임시 기억, 정리하면 공간 확보. 설정 > 저장 공간에서 정리 가능.\textbf{대용량 파일 찾기}: 파일 관리자로 크기별로 정렬하여, 공간을 차지하는 “주범”을 찾아내기.\textbf{잔여물 삭제}: 앱 제거 후 가끔 쓰레기 파일이 남아, 수동으로 정리.\textbf{사진 정리}: 중복 사진, 흐릿한 사진, 불필요한 스크린샷 삭제.\textbf{위챗 정리}: 위챗은 공간의 큰 손이며, 채팅 기록과 캐시를 정기적으로 정리하는 것이 필요.7. 일반적인 파일 형식 완전 분석스마트폰의 다양한 파일 확장자가 머리를 아프게 하는데, 그것을 이해하면 파일을 더 잘 관리할 수 있습니다:형식유형특성권장 열기 방식JPG/JPEG이미지손실 압축, 용량 작음, 호환성 최강. 스마트폰 사진의 기본 형식.갤러리/임의의 이미지 뷰어PNG이미지무손실 압축, 투명 배경 지원. 스크린샷, 아이콘에 자주 사용. 용량은 JPG보다 큼.갤러리/임의의 이미지 뷰어MP3오디오손실 압축, 용량 작음, 호환성 극히 우수. 가장 범용적인 오디오 형식.임의의 음악 플레이어FLAC오디오무손실 압축, 음질이 CD와 일치, 용량은 MP3보다 3-5배 큼. 오디오 애호가 최우선 선택.음악 플레이어/전문가용 재생 앱WAV오디오무압축 원시 오디오, 음질이 최상이지만 용량이 거대. 전문 녹음에 자주 사용.음악 플레이어M4A오디오AAC 인코딩, MP3와 같은 비트레이트에서 음질이 더 좋음. Apple 생태계에서 자주 사용.Apple Music/음악 플레이어MP4영상가장 범용적인 영상 컨테이너 포맷으로, 호환성이 매우 뛰어납니다. 스마트폰에서 녹화된 영상의 기본 포맷입니다.사진 앱/임의의 영상 재생기MKV영상오픈소스 컨테이너 포맷으로, 다중 오디오 트랙과 자막을 지원합니다. 영화 자료에 자주 사용됩니다.VLC/nPlayerDOC/DOCX문서Word 문서. DOCX는 압축 포맷으로 파일 크기가 더 작습니다. DOC은 구버전 포맷입니다.WPS/OfficePPT/PPTX프레젠테이션PowerPoint 프레젠테이션. PPTX는 압축 포맷입니다.WPS/OfficeXLS/XLSX스프레드시트Excel 스프레드시트. XLSX는 압축 포맷입니다.WPS/OfficePDF문서크로스 플랫폼 포맷으로, 레이아웃이 고정되어 편집할 수 없습니다. 계약서, 논문, 전자책에 자주 사용됩니다.Adobe Reader/WPSZIP압축가장 범용적인 압축 포맷으로, 크로스 플랫폼 호환성이 우수합니다.시스템 기본 제공/압축 해제 앱RAR압축ZIP보다 압축률이 높지만, 전용 소프트웨어로 압축 해제해야 합니다.압축 해제AppISO이미지광 디스크 이미지 파일로, 디스크 내용을 완전히 복제합니다. 시스템 설치 패키지에 자주 사용됩니다.압축 해제 앱/가상 광학 드라이브HTML웹페이지웹페이지 파일로, 브라우저로 열어 봅니다. 저장된 웹페이지에 자주 사용되는 포맷입니다.브라우저연습해 보기: “다운로드” 폴더 정리하기파일 관리자를 열고 > “다운로드” 폴더로 이동합니다. 안에 무엇이 있는지 살펴보세요? 필요 없는 것은 삭제하고, 유용한 것은 분류합니다(예: “업무 자료”, “이모티콘” 폴더를 새로 만들어 보관). 그리고 저장 공간도 확인하여 어떤 유형의 파일이 가장 많은 공간을 차지하는지 파악합니다.8. 사진과 파일 시스템 관리: 자산을 깔끔하게 정리하기이전 섹션에서 “폴더 인식”을 배웠지만, 인식만으로는 충분하지 않습니다. 이 섹션에서는\textbf{장기간 실행 가능한 관리 규범을 적극적으로 구축하는 방법}을 다루어, 파일을 언제든 찾을 수 있고, 분실되지 않으며, 도청당하지 않도록 합니다.1. 파일 분류 및 명명 규범혼란의 근본 원인은 규칙이 없기 때문입니다. 규칙을 만드는 것은 사실 매우 간단하며, 핵심은\textbf{지속}에 있습니다.\textbf{“시간+사건” 명명법}: 폴더 이름에 타임스탬프와 사건 설명을 포함하여 나중에 회상과 검색을 용이하게 합니다. 예:\texttt{202606\_여행 영수증},\texttt{2025년 연말 결산},\texttt{2024 아이 생일}. 이렇게 하면 시간대를 확인할 수 있고, 내용도 한눈에 알 수 있습니다.\textbf{3단계 분류 체계}:\textbf{1단계}: 파일 유형별 (문서 / 사진 / 영상 / 오디오)\textbf{2단계}: 사용 시나리오별 (업무 / 생활 / 학습)\textbf{3단계}: 구체적 사건 또는 프로젝트별 (예: “2026Q2\_프로젝트 A”)\textbf{월별 정리 계획}: 매월 고정된 날짜(월급날 또는 매월 1일 추천)에 “다운로드” 디렉토리를 집중적으로 정리합니다:\textbf{유용한 것은 보관, 쓸모없는 것은 삭제, 소중한 것은 별도 저장}. 일 년에 12번의 작은 정리는 일 년치를 모아서 대청소하는 것보다 훨씬 수월합니다.이번 달 작업 목록한가한 주말을 골라 다음 순서대로 진행합니다:파일 관리자에서\texttt{2026Q2\_프로젝트 A}형식의 이번 달 업무 폴더를 생성합니다.“다운로드” 디렉토리에 있는 이번 주 파일을 모두 검토하여, 유용한 것은 남기고, 보관하고, 쓸모없는 것은 삭제합니다.올해의 가족 사진을 위해\texttt{2025가족},\texttt{2026가족}같은 연도별 폴더를 만듭니다.2. 중복 파일 처리 전략스마트폰을 오래 사용하면 “부풀어 오르는데”, 대부분은 중복 파일이 원인입니다.\textbf{72시간 자동 정리}: 위챗, QQ 등 소셜 앱에서 수신하는 임시 파일(이모티콘, 채팅 사진, 작은 영상)은 대부분 한 번 보고 나면 쓸모가 없습니다. [위챗] > [설정] > [일반] > [저장 공간]에서 “자동 정리”를 활성화하고, 만료된 캐시를 72시간으로 설정합니다. 스마트폰이 순식간에 깔끔해지고, 중요한 기록도 잃지 않습니다.\textbf{저장 공간 targeted 정리}: 안드로이드 [설정] > [저장 공간] > [정리 가속], iOS [설정] > [일반] > [iPhone 저장 공간]에서 시스템이 “캐시 / 임시 파일 / 삭제 가능 앱 / 대용량 파일”별로 분류하여 표시합니다. 중점 정리 대상: (1) 캐시 파일; (2) 만료된 설치 패키지(앱 업그레이드 후 남은 구버전 설치 패키지); (3) 중복 파일(같은 이름과 크기의 사진).\textbf{자동 분류 규칙}: 대부분의 브랜드 사진 앱은 인물, 풍경, 문서, 스크린샷 등 지능형 분류를 지원합니다. 활성화 후, 새로 찍은 단체 사진이 자동으로 “인물 > 김씨”로 분류되고, 스캔한 계약서는 자동으로 “문서”로 분류됩니다.\textbf{매달 정리와 함께 사용하면 “물건을 찾을 수 없는” 상황을 크게 줄일 수 있습니다}.3. 사진 보안 처리 규범사진은 스마트폰에서 가장 “민감한” 데이터 유형입니다 - 무심코 찍은 한 장의 사진이 위치, 기기, 시간을 노출할 수 있습니다.\textbf{“이중 보호” 워크플로우 구축}을 권장합니다:신분증 사진/민감한 스크린샷 발송 표준 절차\textbf{가시 정보 마스킹}: 시스템 기본 제공 또는 타사 도구(예: Meitu Xiu Xiu, Snapseed)를 사용하여\textbf{신분증 번호 중간 자리, 차량 번호판, 주소, 이름, 금액}등 핵심 필드를 모자이크 처리합니다.\textbf{메타데이터 제거}: 발송 전 사진 앱에서 “공유” 또는 “내보내기”를 클릭하면 시스템이 위치 정보 유지 여부를 묻습니다.\textbf{“유지하지 않음”}을 선택합니다. 이렇게 하면 촬영 기기 모델, GPS 좌표, 촬영 시간이 사진과 함께 전송되지 않습니다.\textbf{임시 워터마크}(선택): 민감한 스크린샷에 “OO 처리 사용 전용”이라는 텍스트 워터마크를 추가하여 남용을 방지합니다.메타데이터(EXIF) 정리 방법(기본 작업, 심층 방어는 6장 “사진 비식별화 처리” 참조):\textbf{시스템 기본 제공}: iOS [설정] > [프라이버시 및 보안] > [카메라]에서 “정확한 위치”를 끌 수 있습니다. 공유 시 “현재 사진만”을 선택하면\textbf{대부분의 EXIF가 자동으로 제거}됩니다.\textbf{안드로이드}: 대부분의 브랜드 사진 앱은 “공유” 시 “위치 정보 제거 여부”를 묻는 메시지를 표시합니다. 체크하면 됩니다.\textbf{전문 도구}: ExifTool(PC), Metapho(iOS), Photo Metadata Remover(안드로이드)는 모든 메타데이터를 완전히 제거할 수 있습니다.\textbf{장기 보관}이 필요한 민감한 사진(신분증, 은행 카드, 호구부 스캔본, 여권)은\textbf{암호화된 사진 폴더}에 저장할 것을 권장합니다:iPhone: [사진] > 사진 선택 > 숨기기(Face ID 인증 후에만 볼 수 있음);화웨이/샤오미: [사진] > 사진 길게 누르기 > “금고로 이동”;타사 암호화 앱:\textbf{360 보안 사진},\textbf{텐센트 사진 관리자(프라이버시 사진)},\textbf{KeepSafe}등은 별도의 비밀번호로 보호되어 은밀성이 더 좋습니다.}
